\documentclass[11pt, a4paper]{article}

% =====================
% PACKAGES
% =====================
\usepackage{graphicx}
\usepackage{amsmath}
\usepackage{amssymb}
\usepackage{geometry}
\usepackage{booktabs}
\usepackage{hyperref}
\usepackage{listings}
\usepackage{xcolor}
\usepackage{float}
\usepackage{caption}

% =====================
% CONFIG
% =====================
\geometry{margin=1in}
\definecolor{codegray}{rgb}{0.95,0.95,0.95}

\lstset{
    backgroundcolor=\color{codegray},
    basicstyle=\ttfamily\small,
    breaklines=true,
    frame=single,
    columns=fullflexible
}

% =====================
% TITLE
% =====================
\title{\textbf{Project RIFT-MNEMOS:}\\
\large Emergence of Spatially Gated Episodic Memory in a Non-Equilibrium Substrate}

\author{
\textbf{Akshay John} \\
Independent Researcher \\
\texttt{Experiment Series: RIFT-003 / MNEMOS-V}
}

\date{January 2026}

\begin{document}
\maketitle

% =====================
% ABSTRACT
% =====================
\begin{abstract}
[cite_start]Current Artificial Intelligence (AI) architectures, such as Transformers, rely on global optimization (backpropagation) to encode information. [cite: 1] In contrast, biological memory emerges from substrates that must persist under continuous environmental stress. We present \textbf{RIFT} (Resilient Internal Fracture Topology), a neural cellular automaton governed by antagonistic homeostatic regulation. We demonstrate that memory arises not as updated weights, but as \textbf{irreversible topological scarring}. Through a series of failure-driven experiments (MNEMOS I-V), we identify that naive homeostatic regulation leads to "Global Callus" formation, blocking learning. [cite_start]We resolve this by introducing a \textbf{Spatial Novelty Gate}, achieving a memory capacity of 78\% (7/9 sites) in a continuous, non-optimizing substrate. [cite: 5, 22]
\end{abstract}

% =====================
% 1. INTRODUCTION
% =====================
\section{Introduction}
[cite_start]We posit that AGI cannot be built on "Brains in a Jar" (Transformers) which are static and task-dependent. [cite: 8] AGI requires a "Body" that persists in time. Project RIFT investigates the hypothesis: \textit{Can a system encode information purely through physical deformation (scarring) to survive environmental stress?}

% =====================
% 2. ARCHITECTURE
% =====================
\section{System Architecture}
The system is a $48 \times 48 \times 32$ grid governed by three laws:
\begin{enumerate}
    \item \textbf{The Substrate:} A diffusive grid that propagates energy.
    \item \textbf{The Fracture Law:} If local stress $\sigma >$ Elasticity $E$, the grid fractures. [cite_start]This dissipates energy but permanently increases $E$ (Scarring). [cite: 13]
    \item \textbf{The Antagonist Ego:} A global regulator that enforces non-equilibrium. [cite_start]It injects chaos if the system is too stable ($Fracture < 0.005$) and damps energy if the system explodes ($Fracture > 0.15$). [cite: 14]
\end{enumerate}

% =====================
% 3. FAILURE ANALYSIS (The Journey)
% =====================
\section{Methodology: Failure-Driven Refinement}

\subsection{Phase 1: Thermodynamic Instability}
[cite_start]Initial attempts lacked energy dissipation limits, resulting in positive feedback loops and numerical divergence ($Fracture \rightarrow \infty$). [cite: 15]

\subsection{Phase 2: The Global Callus (MNEMOS-II)}
We attempted to store memories by hardening the substrate upon fracture. However, the Antagonist Ego's constant background noise caused the entire grid to harden uniformly.
\textbf{Result:} The system became "numb" to new information. [cite_start]Capacity $\approx 1/9$ sites. [cite: 16]

\subsection{Phase 3: Dual-Timescale Failure (MNEMOS-IV)}
We separated memory into Fast (Skin) and Slow (Bone) layers. [cite_start]This failed because the Fast layer absorbed all stress, effectively "insulating" the Slow layer from ever learning. [cite: 21]

% =====================
% 4. THE SOLUTION (MNEMOS-V)
% =====================
\section{Experiment IV: Spatially Gated Memory (MNEMOS-V)}

\subsection{Hypothesis}
To solve the Signal-to-Noise problem between the Ego (Noise) and the Input (Signal), we introduced a **Spatial Gating Mechanism**. Learning in the permanent memory layer is gated by the "Sharpness" of the stress field:
\[ G(x,y) = |\sigma(x,y) - \mu_{local}| \]
[cite_start]This ensures only structured, high-frequency events (Trauma) cause scarring, while diffuse Ego noise is ignored. [cite: 22]

\subsection{Results}
We subjected the substrate to 9 sequential trauma events at distinct spatial locations.
\begin{lstlisting}
Step 000 | Target X1 | Initial Trauma: 0.2604
...
Step 800 | Target X9 | Initial Trauma: 0.0024
\end{lstlisting}

\textbf{Recall Test:}
\begin{lstlisting}
Recall X1: 0.0023 (Baseline: 0.2604) -> REMEMBERED
Recall X2: 0.0000 (Baseline: 0.0000) -> BLOCKED
...
Recall X9: 0.0000 (Baseline: 0.0024) -> REMEMBERED
\end{lstlisting}

\subsection{Verdict}
\textbf{Sites Retained: 7 / 9}. [cite_start]The system successfully differentiated between homeostatic noise and meaningful environmental stress, encoding the latter as persistent topological scars. [cite: 23]

% =====================
% 5. DISCUSSION & CONCLUSION
% =====================
\section{Conclusion}
RIFT demonstrates that memory is not merely the storage of data, but the \textbf{physical adaptation to trauma}. By implementing spatial gating, we have engineered a non-equilibrium substrate capable of continuous, unsupervised learning. [cite_start]This provides a physical foundation for future AGI architectures that require embodied persistence. [cite: 25]

\end{document}